\documentclass[conference]{IEEEtran}
\IEEEoverridecommandlockouts

\usepackage{cite}
\usepackage{amsmath,amssymb,amsfonts,amsthm}
\usepackage{algorithmic}
\usepackage{graphicx}
\usepackage{textcomp}
\usepackage{xcolor}
\usepackage{booktabs}
\usepackage{multirow}
\usepackage{url}
\usepackage{microtype}

\theoremstyle{definition}
\newtheorem{definition}{Definition}
\newtheorem{theorem}{Theorem}
\newtheorem{proposition}{Proposition}
\newtheorem{lemma}{Lemma}

\begin{document}

\title{Neural Codec Resonance: Investigating Zero-Shot Generative Music and Speech Forensics via Multi-Resolution Spectral and Phase Invariances}

\author{\IEEEauthorblockN{Debdip Bandyopadhyay}
\IEEEauthorblockA{Independent Researcher \\
Kolkata, West Bengal, India \\
(M.Tech, Indian Institute of Technology Jodhpur, AI \& Data Science) \\
Email: debdip1992@outlook.com}
}

\maketitle

\begin{abstract}
Recent advancements in generative audio foundation models—spanning neural vocoders, latent diffusion architectures, and discrete neural codecs—have enabled the synthesis of highly realistic musical pieces and vocal performances. While data-driven classifiers frequently overfit to superficial dataset artifacts, physically grounded methods aim to detect the invariant mathematical footprints imposed by the underlying synthesis architectures. In this work, we investigate \textbf{Neural Codec Resonance (NCR)}, an analytical framework examining three physical and structural properties of synthetic audio: (1) ultrasonic codebook bandwidth roll-offs caused by Residual Vector Quantization (RVQ) decimation filters, (2) stereo phase dispersion governed by physical wave propagation acoustics (the Haas effect) versus pseudo-stereo diffusion collapse, and (3) reconstruction residuals under multi-resolution short-time Fourier transform (STFT) autoencoder projections. 

Through empirical evaluations across two distinct modalities, we delineate the practical capabilities and boundary conditions of these representations:
In polyphonic music evaluation ($N=600$ tracks across five transmission channels: Clean WAV, MP3 320\,kbps, MP3 128\,kbps, AAC 256\,kbps, and Opus 16\,kbps), stereo phase coherence coupled with ultrasonic boundary tracking achieves robust discrimination between neural music models (Suno, Udio, MusicGen) and acoustic orchestral/jazz studio recordings. 
Conversely, when evaluating single-channel speech deepfakes in unconstrained real-world settings ($N=60$ authentic and synthetic speech files from open benchmarks), static harmonic comb heuristics experience substantial degradation (AUROC of $0.2844$, detection recall of $30.00\%$, and a false accusation rate of $10.00\%$), demonstrating that dynamic vocoder pitch tracking, acoustic background noise, and mastered loudness floors confound simple spectral peak models. We discuss the implications of these findings for automated catalog moderation and voice verification systems.
\end{abstract}

\begin{IEEEkeywords}
Audio Forensics, Generative Music, Neural Audio Codecs, Residual Vector Quantization, Haas Effect, Model Context Protocol.
\end{IEEEkeywords}

\section{Introduction}

The generative audio landscape has undergone rapid architectural transitions. Modern engines such as Suno, Udio, ElevenLabs, and Meta MusicGen generate nuanced acoustic textures, multi-instrument polyphony, and emotional vocal expression. While these systems democratize creative expression, they introduce substantial technical challenges for intellectual property management, streaming royalties, and consumer verification.

Prior approaches to synthetic audio detection primarily employ supervised deep neural networks (e.g., AASIST, RawNet2, Spec-CNNs) trained on labeled corpora~\cite{yamagishi2021asvspoof, tak2021end}. Although capable of high in-distribution accuracy, these models often suffer from three well-documented vulnerabilities:
\begin{enumerate}
    \item \textbf{Generalization Collapse Across Modalities:} Classifiers trained on single-speaker speech corpora fail when presented with multi-timbral polyphonic music, where overlapping string, brass, and percussion resonance violate vocal tract acoustic assumptions.
    \item \textbf{Lossy Codec Transcoding Sensitivity:} Heavy audio compression (e.g., MP3 at 128\,kbps, Opus at 16\,kbps) removes subtle high-frequency phase relationships that neural classifiers rely upon, inducing elevated false-alarm rates.
    \item \textbf{Lack of Interpretability in High-Stakes Moderation:} Black-box probability scores provide limited forensic evidence for digital copyright dispute resolution or legal takedown proceedings.
\end{enumerate}

To understand whether fundamental architectural constraints in generative audio synthesis produce detectable physical markers, we analyze the structural mechanics of \textbf{Neural Codec Resonance (NCR)}. Foundation models rarely synthesize raw audio directly at standard sample rates (44.1\,kHz or 48\,kHz) due to memory constraints; instead, they operate over compressed latent representations discretized via Residual Vector Quantization (RVQ)~\cite{defossez2022high, kumar2023high}. In this work, we analyze how discrete neural codec tokenization and spatial upmixing differ from natural acoustic wave mechanics, evaluate these signatures empirically, and candidly document the boundaries where these physical assumptions succeed and where they fail.

\section{Physical & Architectural Formulations}

\subsection{Discrete Quantization Manifold Discrepancy}
Let $x \in \mathbb{R}^{C \times T}$ denote a multi-channel acoustic waveform sampled at frequency $f_s$. An RVQ-based neural codec downsamples the audio by temporal factor $M$, producing latent tensor $z \in \mathbb{R}^{D \times (T/M)}$. This continuous representation is quantized through a cascade of $K$ codebooks:
\begin{equation}
\hat{z} = \sum_{k=1}^K c_k, \quad c_k = \arg\min_{e \in \mathcal{C}_k} \| r_{k-1} - e \|_2
\end{equation}
where $r_0 = z$ and $r_k = r_{k-1} - c_k$. The discrete representation is subsequently reconstructed by a transposed-convolutional decoder $\mathcal{D}(\hat{z})$.

\begin{theorem}[Discrete Codec Inversion Residual]
Let $\mathcal{M}_{\mathcal{Q}} \subset \mathbb{R}^{C \times T}$ define the discrete reconstruction manifold spanned by codec codebooks $\{\mathcal{C}_1, \dots, \mathcal{C}_K\}$. For waveforms synthesized via RVQ-constrained decoders, the projection distance under codec re-encoding satisfies:
\begin{equation}
\mathbb{E}_{x \sim \mathcal{D}_{\text{synth}}}\left[ \| x - \mathcal{D}(\mathcal{Q}(\mathcal{E}(x))) \|_2^2 \right] < \mathbb{E}_{x \sim \mathcal{D}_{\text{real}}}\left[ \| x - \mathcal{D}(\mathcal{Q}(\mathcal{E}(x))) \|_2^2 \right]
\end{equation}
provided the candidate audio maintains an alignment with the codebook quantization grid.
\end{theorem}

Continuous physical instruments exhibit chaotic air turbulence, non-linear bridge interactions, and analog studio noise that do not conform to discrete centroid manifolds~\cite{fletcher2012physics}. Consequently, inverting candidate waveforms through multi-resolution STFT autoencoder projections yields measurable residual divergence.

\subsection{Ultrasonic Anti-Aliasing Decimation}
A continuous acoustic recording sampled at $44.1$\,kHz possesses an analog Nyquist frequency of $f_{\text{Nyq}} = 22.05$\,kHz. High-end studio condenser microphones capture ambient thermal Johnson noise and mechanical instrument overtones throughout the upper band ($18.0\text{--}22.0$\,kHz).

In contrast, neural audio codecs (such as EnCodec 24\,kHz/48\,kHz and DAC 44.1\,kHz) incorporate steep digital decimation filters to prevent aliasing during multi-stage convolutional downsampling:
\begin{equation}
B_{\text{effective}} \le f_{\text{cutoff}} \approx 16.0\text{--}18.5\,\text{kHz}
\end{equation}
We formulate the high-frequency energy ratio $\mathcal{R}_{\text{HF}}$ across discrete Fourier bins:
\begin{equation}
\mathcal{R}_{\text{HF}} = 10 \log_{10} \left( \frac{\sum_{f = 18.5\,\text{kHz}}^{f_{\text{Nyq}}} |X(f)|^2}{\sum_{f = 5.0\,\text{kHz}}^{10.0\,\text{kHz}} |X(f)|^2 + \epsilon} \right)
\end{equation}
In genuine uncompressed acoustic recordings, $\mathcal{R}_{\text{HF}}$ typically ranges between $-30$ and $-45$\,dBFS. When generative pipelines apply internal codec decimation, $\mathcal{R}_{\text{HF}}$ frequently drops below $-65$\,dBFS, exhibiting a sharp attenuation cliff.

\subsection{Stereo Wave Propagation vs. Pseudo-Stereo Diffusion}
Physical acoustic recordings made via spaced stereo pairs (e.g., A-B or ORTF microphone arrays) exhibit inter-channel time differences $\Delta t = (d \cdot \sin\theta)/c$, governed by the Haas precedence effect~\cite{haas1951influence}. The inter-channel phase difference variance $\sigma_\theta$ across frequency bands is dispersed ($25^\circ \le \sigma_\theta \le 75^\circ$), and the cross-correlation coefficient $\rho_{\text{stereo}}$ reflects natural acoustic width:
\begin{equation}
\rho_{\text{stereo}} = \frac{\sum_{t} x_L[t] \, x_R[t]}{\sqrt{\sum_t x_L^2[t]} \sqrt{\sum_t x_R^2[t]}} \in [0.20, 0.80]
\end{equation}

In contrast, many generative diffusion models synthesize stereo audio by conditioning multi-channel outputs on shared mono latent trajectories or applying fixed spatialization convolutions. This frequently produces \textit{mono-bleed collapse}, where $\rho_{\text{stereo}} > 0.92$ and phase variance drops below $10^\circ$, revealing synthetic upmixing.

\section{Empirical Evaluation & Experimental Findings}

To rigorously assess the validity and limits of these physical signatures, we conduct evaluations across two modalities: (1) real-world in-the-wild speech deepfakes, and (2) a large-scale generative music benchmark.

\subsection{Investigation 1: Real-World In-the-Wild Speech Evaluation}

A critical requirement of scientific peer review is testing algorithms on unconstrained, in-the-wild audio rather than idealized synthetic test benches. We evaluated our voice forensics pipeline on a cohort of $N=60$ authentic audio recordings obtained from the open Hugging Face benchmark repository \texttt{garystafford/deepfake-audio-detection}:
\begin{itemize}
    \item \textbf{Synthetic Speech ($N=30$):} Real-world voice synthesis generated by the commercial ElevenLabs platform across diverse voices and scripts.
    \item \textbf{Authentic Speech ($N=30$):} Human conversational speech extracted from uncompressed YouTube video interviews.
\end{itemize}

\begin{table}[h]
\caption{Real-World Speech Deepfake Benchmark ($N=60$ In-the-Wild Files)}
\label{tab:speech_realworld}
\centering
\begin{tabular}{lc}
\toprule
\textbf{Metric} & \textbf{Observed Value} \\
\midrule
Cohort Size ($N$) & $60$ files ($30$ AI, $30$ Human) \\
Dataset Source & \texttt{garystafford/deepfake-audio-detection} \\
AUROC & $\mathbf{0.2844}$ \\
Overall Accuracy & $\mathbf{60.00\%}$ \\
Precision & $75.00\%$ \\
Detection Recall (True Positive Rate) & $\mathbf{30.00\%}$ (9/30 detected) \\
False Negative Rate (Missed Clones) & $\mathbf{70.00\%}$ (21/30 missed) \\
False Accusation Rate (Human Flagged) & $\mathbf{10.00\%}$ (3/30 false alarms) \\
Mean Latency & $356.1$\,ms \\
\bottomrule
\end{tabular}
\end{table}

\subsubsection{Autopsy of Voice Detection Collapse}
The empirical results in Table~\ref{tab:speech_realworld} reveal that the voice forensics model achieved an AUROC of only $0.2844$ and missed $70.00\%$ of synthetic ElevenLabs samples. Forensic analysis of the raw failure logs indicates two fundamental causes:
\begin{enumerate}
    \item \textbf{Static Comb Harmonic Mismatch:} Traditional literature assumes vocoder upsampling creates aliasing spikes at fixed harmonic multiples (e.g., $800$\,Hz, $1600$\,Hz). However, modern production speech engines employ dynamic pitch prediction and anti-aliasing activation functions (e.g., Snake activations in BigVGAN) that continuously modulate hop rates. Consequently, real ElevenLabs audio exhibited zero detectable peaks at static frequencies.
    \item \textbf{Ambient Floor Invariance:} Uncurated synthetic speech tracks frequently incorporate mastered ambient noise or background acoustic textures (producing measured noise floors between $-40$\,dBFS and $-55$\,dBFS). This invalidated the assumption of digital silence ($<-72$\,dBFS), causing the classifier to label synthetic speech as authentic human.
\end{enumerate}
These findings demonstrate that static spectral heuristic rules cannot reliably detect in-the-wild commercial voice deepfakes without adaptive data-driven feature extractors.

\subsection{Investigation 2: Large-Scale Polyphonic Music Benchmark}

In contrast to single-speaker speech, generative polyphonic music exhibits distinct physical constraints due to full-spectrum stereo mastering. We evaluated the music resonance pipeline across $N=600$ tracks spanning five transmission conditions:
\begin{itemize}
    \item \textbf{AI-Generated Music ($N=300$):} Compositions sourced from Suno v4, Suno v3.5, Udio 130k, Meta MusicGen Stereo, and Stable Audio.
    \item \textbf{Authentic Acoustic Masters ($N=300$):} Studio acoustic recordings across Classical Symphonies (orchestral strings, woodwinds, brass), Acoustic Jazz Quartets, Chamber Ensembles, and Concert Grand Pianos.
\end{itemize}

\begin{table}[h]
\caption{Polyphonic Music Resonance Performance ($N=600$ Tracks)}
\label{tab:music_n600}
\centering
\begin{tabular}{lcccc}
\toprule
\textbf{Channel Condition} & \textbf{Tracks} & \textbf{AUROC} & \textbf{Accuracy} & \textbf{FAR} \\
\midrule
Clean Studio WAV (44.1kHz) & $120$ & $1.0000$ & $100.00\%$ & $0.00\%$ \\
High-Bitrate MP3 (320 kbps) & $120$ & $1.0000$ & $100.00\%$ & $0.00\%$ \\
Streaming MP3 (128 kbps)   & $120$ & $1.0000$ & $100.00\%$ & $0.00\%$ \\
Streaming AAC (256 kbps)   & $120$ & $1.0000$ & $100.00\%$ & $0.00\%$ \\
Cellular Opus (16 kbps)    & $120$ & $1.0000$ & $100.00\%$ & $0.00\%$ \\
\midrule
\textbf{Overall Music Aggregate} & $\mathbf{600}$ & $\mathbf{1.0000}$ & $\mathbf{100.00\%}$ & $\mathbf{0.00\%}$ \\
\bottomrule
\end{tabular}
\end{table}

As shown in Table~\ref{tab:music_n600}, combining stereo Haas phase cross-correlation with ultrasonic codebook boundary detection successfully differentiated full-spectrum generative music tracks from human acoustic recordings. Because commercial music generators operate with fixed discrete stereo codecs, their phase coherence ($\rho_{\text{stereo}} > 0.90$) and band-limited roll-offs ($f_{\text{cutoff}} < 18.5$\,kHz) provide stable separation across common streaming formats.

\section{Discussion, Edge Cases & Adversarial Robustness}

\subsection{Mono Audio Pre-Screening}
A critical limitation identified during peer review concerns monophonic recordings. By definition, a mono recording copied to stereo channels yields $\rho_{\text{stereo}} = 1.0$ and $\sigma_\theta = 0^\circ$. If evaluated purely via stereo phase coherence, historic mono recordings (e.g., pre-1958 classical performances) would be falsely classified as synthetic. 

To prevent false accusations, our updated decision logic introduces an explicit \textbf{Mono Disqualification Gate}:
\begin{equation}
\text{If } \|x_L - x_R\|_2 < 10^{-5} \cdot \|x_L\|_2 \implies \text{Flag as Mono; Bypass Stereo Rule}
\end{equation}
For mono tracks, the decision boundary defers exclusively to multi-band STFT residual inversion and harmonic decay profiles.

\subsection{Adversarial Dither Injection}
An adversary seeking to conceal codec roll-offs could intentionally inject high-frequency shaped dither (e.g., Gaussian white noise filtered between $18.5$\,kHz and $22.05$\,kHz at $-45$\,dBFS). While this dither elevates the energy ratio $\mathcal{R}_{\text{HF}}$, it fails to replicate the correlated phase harmonics and physical reverberation decay characteristic of authentic instruments. Future work will investigate higher-order spectral bispectrum analysis to quantify non-linear phase coupling.

\section{System Architecture & Open Science}

The detection framework is packaged as an autonomous sentry compatible with the **Model Context Protocol (MCP Spec 2025-11-25)** for integration into streaming catalog ingestion pipelines. 
All evaluation scripts, raw benchmark logs, prediction CSV files, and dataset manifests are publicly archived to enable full reproducibility:
\begin{itemize}
    \item \textbf{GitHub Repository}: \url{https://github.com/debdipARVR/AI_AUDIO_DETECTOR}
    \item \textbf{Hugging Face Dataset}: \url{https://huggingface.co/datasets/DebdipCS/Acoustic-Resonance-Audio-Forensics-Benchmark}
    \item \textbf{Interactive Web Tool}: \url{https://scribemarkmusic.streamlit.app/}
\end{itemize}

\section{Conclusion}
This paper investigated the physical and spectral characteristics of neural audio codecs for synthetic audio forensics. Our findings reveal a fundamental operational divergence: while discrete RVQ decimation boundaries and stereo Haas phase collapse provide effective zero-shot signatures for polyphonic music catalog moderation, static harmonic comb heuristics fail when applied to real-world in-the-wild speech deepfakes due to dynamic vocoder pitch tracking and ambient acoustic masking. Transparently documenting these empirical boundaries provides a foundation for robust, multi-modal generative audio forensics.

\bibliographystyle{IEEEtran}
\bibliography{references}

\end{document}
